\medskip\noindent
\textbf{All stable degree-eight components and their actual maps.}
Work two-locally, and take a finite stable approximation to
$\Sigma^\infty F$ through dimension ten before cancelling odd-unit
cell pairs. Retain its nine-skeleton, with lower skeleta $L_i$.
This preserves the required stable group and removes extraneous
nine-cells killed by unit ten-cell incidences. The lower model is
the one in D.2.18--D.2.19:
\[
 L_6=C\vee M(\ZZ/4,5)\vee S^6_{\rho_0},
 \qquad
 C=\cofib\bigl((2,\eta_4):S^4\vee S^5\to S^4\bigr).
 \tag{D.2.20t1}
\]
All groups of the $L_i$ in this calculation are stable groups.
The quotient $C\to S^5\vee S^6$ injects $\pi_8C$ into
$\ZZ/8\{\nu_5\}\oplus\ZZ/2\{\eta_6^2\}$.
Indeed the preceding term is $\pi_4^S=0$, and the following
map is $(x,y)\mapsto2x+4y$ in $\ZZ/8$.
Its kernel is cyclic of order four. Choose its generator $\omega$
so that
\[
 \pi_8C=\ZZ/4\{\omega\},\qquad
 q_*\omega=(2\nu_5,\eta_6^2),\qquad
 \delta\eta_6^2=2\omega.
 \tag{D.2.20t2}
\]
The last equality follows from $q_*\delta=(\eta_5,0)$ and
$\eta^3=4\nu$ two-locally. The stable stems and Moore
coextension formulas used here are the opened formulas in
\cite[\S2, (2.1)--(2.2), p.~149]{I_Oka1986}.

For the other summands, the Moore exact sequence and the
order-two lift $b_0$ of D.2.19 give
\[
\begin{split}
 \pi_8M(\ZZ/4,5)
   &=\ZZ/4\{n_0=\iota_5\nu_5\}
       \oplus\ZZ/2\{b_0\eta_7\},\\
 \pi_8S^6_{\rho_0}&=\ZZ/2\{c_0\eta_7\},
       \qquad c_0=\rho_0\eta_6.
\end{split}
 \tag{D.2.20t3}
\]
Here $\iota_5$ is the bottom Moore inclusion, whose
homology class is $u_5$.
The top pinch sends $b_0\eta_7$ to $\eta_6^2$;
it has order two because $b_0$ does.
Also $t\eta_6^2=i\eta^3=0$ in the mod-four Moore summand.

The seven-cell attachments
\[
 \gamma+t,\quad \gamma+2\rho_0+\delta,\quad\delta,\quad0
\]
have $\eta^2$-images $0,2\omega,2\omega,0$.
Their maps on the preceding $\eta$-groups have kernel
generated by $\eta$ on the cell combinations $P+Q$ and $R$.
Thus the cofiber sequence gives
\[
\begin{split}
 0\longrightarrow&
 \ZZ/2\{\overline\omega\}\oplus\ZZ/4\{n_0\}
 \oplus\ZZ/2\{b_0\eta_7,c_0\eta_7\}\\
 \longrightarrow&\ \pi_8L_7
 \longrightarrow (\ZZ/2)^2\longrightarrow0.
\end{split}
 \tag{D.2.20t4}
\]
The $R$ class has the actual order-two lift
$\ell_0=\lambda\eta_7$, since that seven-cell is freely wedged.
For the other lift, change the seven-cell basis from $P,Q$ to
$P+Q,Q$. The new $P+Q$ attachment is $\gamma+2\rho_0$.
Choose a lift $\vartheta$ of its top $\eta$ in the subcomplex
\[
 D=(C\vee S^6_{\rho_0})
          \mathop{\cup}_{\gamma+2\rho_0}e^7_{P+Q}.
 \tag{D.2.20t5}
\]
It exists because $\gamma\eta_6=0$ and $2\rho_0\eta_6=0$.
The element $2\vartheta$ comes from $\pi_8(C\vee S^6)$;
in $\pi_8L_7$ it therefore belongs to the subgroup generated
by $\overline\omega$ and $c_0\eta_7$. In particular it has
no $n_0$ component. This specified representative excludes
an ambiguity that an abstract group extension alone would leave.

\medskip\noindent
\textbf{The precise $\kappa\eta$ relation.}
Set $C'=C\cup_\delta e^7_Q$. Equation
\textup{(D.2.20t2)} and $\delta\eta_6\ne0$ give
\[
 \pi_8C'=\ZZ/2\{\overline\omega\}.
\]
Choose $\kappa'\in\pi_7C'$ with top-cell degree two on $Q$.
In its cone the new degree-eight class has Bockstein $Q$.
Since $\operatorname{Sq}_*^2Q=z_5$ and
$\beta z_5=A_4$, the Adem identities force its
second-square vector to be $e_6$, the only six-class in $C'$.
The quotient of $C'$ by its bottom sphere is
\[
 C'/S^4\simeq C\eta_{5}\vee S^6_{e_6},
\]
where $C\eta_5$ has cells in dimensions five and seven.
The $S^6$ component of $\kappa'$ is therefore $\eta_6$.
Consequently $\kappa'\eta_7$ has a nonzero $\eta_6^2$
component in this quotient, and hence equals
$\overline\omega$.

The normalized $\kappa$ of D.2.19 has second-square vector
$d_6$. Both $\kappa$ and $\kappa'$ project to twice the
$Q$ cell. Their difference in $\pi_7L_7$ is therefore
torsion. The complete detectors of D.2.19 show that this
difference is $b_0$ plus possibly $v_0$: the required
second-square difference is $d_6+e_6$, and no $c_0$ term
is possible. As $v_0\eta_7=i\nu\eta=0$, we obtain
\[
 \kappa\eta_7=\overline\omega+b_0\eta_7.
 \tag{D.2.20t6}
\]
This calculation uses the fixed quotient normalization of
$\gamma$ and the fixed second-square normalization of
$\kappa$ from the preceding claims.

\medskip\noindent
\textbf{The entire group before the nine-cells.}
The eight-cell attachments of D.2.19, ordered as
$W_8,C_8,L_8,M_8,J_1,J_2,J_3$, are
\[
 0,\quad v_0,\quad\kappa,\quad
 \kappa+b_0+c_0,\quad b_0,\quad0,\quad2\lambda+b_0.
\]
Their $\eta$-images generate exactly
$\langle\kappa\eta_7,b_0\eta_7,c_0\eta_7\rangle$.
By \textup{(D.2.20t6)}, these kill
$\overline\omega,b_0\eta_7,c_0\eta_7$ in
\textup{(D.2.20t4)}. They also kill $2\vartheta$, by
the representative in \textup{(D.2.20t5)}, and impose
no relation on $n_0$. The resulting torsion subgroup is exactly
\[
 T_8=\ZZ/4\{n_0\}\oplus
        \ZZ/2\{\vartheta,\ell_0\}.
 \tag{D.2.20t7}
\]
The kernel on the following free degree-seven cell groups
is freely generated by
\[
 W_8,\quad J_2,\quad2C_8,\quad2J_1,\quad2(M_8-L_8).
\]
The extension by this free group splits. More invariantly,
write $D_8=M_8-L_8$ in integral cellular coordinates.
The integral Hurewicz map of $L_8$ has kernel precisely
$T_8$ and image
\[
 \Lambda=\ZZ_{(2)}\{W_8,J_2\}
       \oplus2\ZZ_{(2)}\{C_8,J_1,D_8\}
 \subseteq H_8(L_8;\ZZ_{(2)}).
 \tag{D.2.20t8}
\]
In particular the free lift of $2C_8$ will be denoted $a_0$.

For an attaching map with zero integral boundary, the complete
torsion detectors needed here are
\[
\begin{array}{c|c|c}
 &\operatorname{Sq}_*^2&\operatorname{Sq}_*^4\\ \hline
 n_0&0&u_5\\
 \vartheta&P+Q&z_5\\
 \ell_0&R&0.
\end{array}
 \tag{D.2.20t9}
\]
The $n_0$ row is the Hopf-invariant-one cone on the
Moore bottom, and the $\ell_0$ row is the cone of $\eta$
on the $R$ seven-sphere. For $\vartheta$, its top quotient
is $\eta$ on $P+Q$, giving its second square. Its chosen
representative in \textup{(D.2.20t5)} has no $u_5$ class,
so its fourth-square $u_5$ coefficient is zero. Finally,
dualize the Adem identity
$\operatorname{Sq}^2\operatorname{Sq}^3
=\operatorname{Sq}^5+\operatorname{Sq}^4\operatorname{Sq}^1$.
The top Bockstein is zero, while
$\operatorname{Sq}_*^3(P+Q)=A_4$ and $\beta z_5=A_4$.
Its fourth-square $z_5$ coefficient must therefore be one.
For the order-four class $n_0$, these primary detectors
see only the parity of its coefficient; we retain that
limitation in the next step.

\medskip\noindent
\textbf{Three zero-boundary relations kill the entire torsion subgroup.}
The genuine integral cycles $N_9$ and $K_2+K_3+K_4$ constructed
above can be chosen as zero-boundary nine-cells. By their
operation vectors and \textup{(D.2.20t9)}, their attachments
have the respective forms
\[
 \ell_0+(2r+1)n_0,\qquad \ell_0+2s n_0.
 \tag{D.2.20t10}
\]
Neither has a $\vartheta$ component. Subtracting the two
relations gives an odd multiple of $n_0$, so they kill
the entire order-four class $n_0$, and then $\ell_0$.
No value of an undetected even coefficient has been assumed.

The additional cycle $Z_0=P_m+Z_A+P_n$ has zero first and
second Bocksteins by \textup{(D.2.20h9)}. The integral
normalization explained above makes it a zero-boundary
nine-cell. Its operation vector is $(P+Q,z_5)$, so its
attachment is $\vartheta$ plus an even multiple of $n_0$,
with no $\ell_0$ term. It kills $\vartheta$ after the
preceding two relations. Thus all of $T_8$ is killed.
Both second Bocksteins in \textup{(D.2.20h9)} are essential
here; omitting the one on $P_m+Z_A$ would leave an incorrect
order-four stable summand.

\medskip\noindent
\textbf{The remaining quotient is the integral Hurewicz image.}
Let $\mathcal B_9$ be the image of the integral ninth
cellular boundary in $H_8L_8$. Every such boundary belongs
to the lattice $\Lambda$. Since the whole Hurewicz kernel
$T_8$ has already been killed by zero-boundary cells,
the complete stable homotopy quotient is exactly
\[
 \pi_8\Sigma^\infty F=\Lambda/\mathcal B_9.
 \tag{D.2.20t11}
\]
Two attachments with the same integral boundary differ
by $T_8$, so no free-lift extension remains in this formula.
In the integral homology calculation, the four
first-Bockstein images supply order-two lifts of
$W_8,J_1,J_2,D_8$, and $C_8$ is the order-four generator.
Consequently \textup{(D.2.20t8)} identifies the quotient
with the following subgroup of $H_8F$:
\[
 \pi_8\Sigma^\infty F
 \xrightarrow[\cong]{\ h_{\ZZ}\ }
 \langle2C_8,W_8,J_2\rangle
       \subseteq\ZZ/4\{C_8\}\oplus
           \ZZ/2\{W_8,J_1,J_2,D_8\}.
 \tag{D.2.20t12}
\]
Changing an order-two lift by $2C_8$ does not change
this subgroup. It is $(\ZZ/2)^3$. In particular, if
$a_0$ denotes the class with integral Hurewicz image
$2C_8$, then
\[
 \ker\bigl(h_{\FF_2}:\pi_8\Sigma^\infty F
                     \to H_8(F;\FF_2)\bigr)
       =\ZZ/2\{a_0\}.
 \tag{D.2.20t13}
\]
All nine-cell attachments have been accounted for:
the zero-boundary ones killed $T_8$, and their integral
boundary subgroup gives exactly \textup{(D.2.20t11)}.
Higher cells cannot change this stable group.

\medskip\noindent
\textbf{The ordinary Hurewicz map is injective.}
Write $F=\Omega^2 Z$. The double suspension--loop unit
\[
 F\longrightarrow\Omega^2\Sigma^2F
\]
has a homotopy retraction by the adjunction triangle
identity. Since $\Sigma^2F$ is five-connected, its
tenth homotopy group is stable. Therefore stabilization
gives a split injection
\[
 \pi_8F\lhook\joinrel\longrightarrow\pi_8\Sigma^\infty F.
 \tag{D.2.20t14}
\]
In particular the stable zeros $n_0=0$ and $\ell_0=0$
already prove the actual zeros in \textup{(D.2.20h2)}.

The remaining stable mod-two-Hurewicz-kernel element
$a_0$ is not an actual sphere map. Its integral
Hurewicz image is $2C_8$, with $C_8$ of order four.
In its stable cone a new degree-nine mod-two class $x_9$
satisfies
\[
 \beta x_9=C_8\pmod{\operatorname{im}\beta|_{H_9F}}.
\]
The old first-Bockstein images are
$W_8,J_1,J_2,L_8+M_8$; all have zero fourth square.
Since $\operatorname{Sq}_*^4C_8=A_4$, the Adem identity
$\operatorname{Sq}^1\operatorname{Sq}^4
=\operatorname{Sq}^5$ gives
\[
 \operatorname{Sq}_*^5x_9
   =\operatorname{Sq}_*^4\beta x_9=A_4\ne0.
 \tag{D.2.20t15}
\]
If $a_0$ came from an actual map $S^8\to F$, its
ordinary space cone would have this operation.
But $\operatorname{Sq}^5$ on degree-four cohomology
is zero by instability. This contradiction excludes
$a_0$ from the image of \textup{(D.2.20t14)}.
By \textup{(D.2.20t13)}, that image meets the stable
mod-two Hurewicz kernel only in zero. Hence
$\pi_8F\to H_8(F;\FF_2)$ is injective, as claimed.
No choice of a stable quadratic complement and no
identification of stable and ordinary eighth homotopy
groups has been used.
